\subsection{A degree-eight replacement compatible with the infinite class-two tower}

\begin{statement}{D.33}{Proved}{A class-four two-primary kernel with third subgroup of rank 9105}
This addresses the genuine realization question in Part D(e),
manuscript \S51.16, folio 734 verso, scan 174, and the two-CM-coordinate
follow-up to \S51.15, folios 725--727, scans 156--160.
Use $K,N,Q,\mathcal P_2$ and $H^{(k)}$ from D.29--D.32, and set
\[
 T_3=B/D_9(K),\qquad U_8=\gamma_3N\,D_9(K)/D_9(K),\qquad
 C_{***}=B/(\gamma_3N\,D_9(K)).                         \tag{D.33.1}
\]
The elementary abelian group $U_8$ has dimension $9105$, with successive
degree-six, degree-seven and degree-eight dimensions $240,1620,7245$.
The image of $\gamma_4N$ in $T_3$ has dimension $1620$, and the image of
$\gamma_5N$ is trivial. There is an explicit new peripheral automorphism of
$H^{(2)}\times_{C_{***}}T_3$, with multiplier $1$, retaining the homology
and extra CM pair of D.27. The exact orders are
\[
 |T_3|=576\cdot2^{11563},\qquad
 |C_{***}|=576\cdot2^{2458},\qquad
 |H^{(2)}\times_{C_{***}}T_3|=576\cdot2^{18708}.        \tag{D.33.2}
\]
Its complete $12,921,196$-letter cusp word is the identity in both full
finite models. The independently collapsed words have $2,128,528$ letters
in the rank-four $K$ basis and $836,592$ letters in the rank-97 $N$ basis.

The $T_3$ component of this action extends to the precise quotient
\[
 \mathcal S_2^{(8)}=
 \mathcal P_2/\bigl(\overline{\gamma_3N^{(2)}}\cap D_9(K^{(2)})\bigr)
 \simeq \mathcal P_2^{[2]}\times_{C_{***}}T_3.          \tag{D.33.3}
\]
The first lookahead may replace the initial $H^{(2)}$ central coordinates;
the entire $T_3$ action is retained.
The kernel over $Q$ has class exactly four, with
\[
 \gamma_3=C_2^{9105},\qquad \gamma_4=C_2^{1620},\qquad \gamma_5=1.
                                                               \tag{D.33.4}
\]
The compatible multiplier has two-adic component in $1+32\mathbb Z_2$;
all odd-prime components can be $1$. An actual first continuation to
the rank-97 class-two model over $\mathbb Z/8$ has multiplier $1$ and
an $18,077,708$-letter cusp word. It retains the entire $T_3$ action.
The resulting action combines with all odd relative pro-prime
completions and D.28's mixed quotient $J$ along their specified common
quotient $H\times_QG_3$.

There is a canonical surjection
\[
 \mathcal S_2^{(8)}\longrightarrow\mathcal S_2^{(7)}
 \quad\hbox{with kernel }C_2^{7245}.                  \tag{D.33.5}
\]
The induced action on its target replaces the recorded D.31 action:
the images of $s$ first differ in degree five of the $K$ Magnus expansion.
Thus D.32's obstruction to that recorded action is retained. The new
action supplies compatible nonarithmetical shadows on characteristic
refinements retaining the $G_6$ and CM data. Attainability of the same
CM pair by a genuine full profinite peripheral automorphism remains
undecided. D.45 later proves that this particular recorded degree-eight
action has no full relative pro-two peripheral lift.
\end{statement}
\begin{proof}
We first determine the common quotient using actual group generators.
The images of $D_j(K)$ in $N/\Phi(N)$, for $3\leq j\leq8$, have dimensions
$88,68,43,19,4,0$. Choose independent images starting with the deepest
layer, using the retained Hall-power words, and then adjoin the nine
degree-two relation words. This gives the following basis of the
97-dimensional Frattini quotient:
\[
\begin{array}{c|rrrrrr}
 \text{least }K\text{ degree}&2&3&4&5&6&7\\\hline
 \text{number of generators}&9&20&25&24&15&4.
\end{array}                                                     \tag{D.33.6}
\]
All generators are explicit words in $B$. They are a basis of $N^{(2)}$:
their images generate its Frattini quotient, so the endomorphism of the
free pro-two group of rank 97 sending its original basis to these words
is surjective, hence an automorphism. The change of basis and its
homology matrix are retained in the certificate. This argument makes
no assertion that the weighted basis is an abstract discrete free basis.

Collection in these actual generators bounds the image of $\gamma_3N$
modulo $D_9(K)$. Commutators of length at least five have degree at least
ten. The contributing basic commutators of length three and four have
weighted counts
\[
\begin{array}{c|r}
 \text{degree and pattern}&\text{number}\\\hline
 6:(2,2,2)&(9^3-9)/3=240\\
 7:(2,2,3)&9^2\cdot20=1620\\
 8:(2,2,4)&9^2\cdot25=2025\\
 8:(2,3,3)&9\cdot20^2=3600\\
 8:(2,2,2,2)&(9^4-9^2)/4=1620.
\end{array}                                                     \tag{D.33.7}
\]
Products are additive on this subgroup and squares vanish, since its
least degree is six. These counts therefore give the upper bound
$240+1620+7245=9105$. The actual words from D.31 provide the first
1860 independent vectors, retaining their degree-eight terms. In degree
eight the words $[r_i,[r_j,z]]$, with $z$ ranging over all 70 restricted
degree-four Hall-power words of $K$, have rank 3645. Adjoining
$[z_i,[r_j,z_k]]$, with $z_i,z_k$ ranging over the 20 degree-three
Lyndon words, raises the rank to 7245. Their leading terms are realized
by precisely these commutator words. Thus the lower bound equals the
collection bound. Fourfold commutators in the nine $r_i$ have rank
1620. This proves (D.33.4) for the image in $T_3$, including the exact
fourth subgroup rather than only its ambient Lie span.

The $N$ kernel of $T_3$ has order $2^{11563}$ by D.32. Quotienting by
$U_8$ gives $2^{2458}$. There is also an explicit class-two description
of this common quotient. In the weighted basis (D.33.6), retain each
generator and each central square or commutator until its doubled
weighted degree reaches nine. Of the 4753 square/commutator coordinates,
45 have weight four, 2271 have weights five through eight, and 2437
have weight at least nine. Consequently the normal-form order is
\[
             2^{97+2316+45}=2^{2458}.                  \tag{D.33.8}
\]
The map from this presentation onto the actual quotient has the same
order on both sides, so it is an isomorphism. In particular the exact
kernel of $H^{(2)}\to C_{***}$ is
$(\mathbb Z/4)^{2437}\times(\mathbb Z/2)^{2271}$: take the coordinates
of weight at least nine, and the squares of those of weights five
through eight. They are actual words in $D_9(K)$.
Furthermore
\[
 (\gamma_3N[N,N]^4N^8)D_9(K)=\gamma_3N D_9(K),        \tag{D.33.9}
\]
since the two exponent factors lie in $D_{16}(K)$. This is the exact
common-quotient identity. It proves the fibre-product assertion,
surjectivity of $B$ onto that product, and all orders in (D.33.2).

We next specify a correcting group whose law is genuinely abelian in
both targets. Keep the nine original $N$ generators with independent
degree-two terms. Replace each of the other 88 generators by itself
times the inverse of the product of those nine that cancels its
degree-two term. These triangular Nielsen operations give an actual
discrete free basis. Its 4708 squares and commutators involving at least
one of the latter 88 generators have degree at least five. The remaining
45 have degree four; their squares have degree eight. Adjoin the 9105
actual generators of $U_8$. In the joint quotient these words generate
\[
       \mathcal A\simeq(\mathbb Z/4)^{4708}
                         \times(\mathbb Z/2)^{9150}.   \tag{D.33.10}
\]
They commute in $T_3$ because each has degree at least five, and commute
in $H^{(2)}$ because their images lie in its abelian subgroup $C_1$.
The 4708 high words and the 45 twice-low words have independent images
in $C_1$, of order $4$ and $2$ respectively. Their kernel in the joint
group is exactly $U_8$, proving (D.33.10). Equivalently this is the full
image of $D_5(K)\cap\Phi(N)$. Indeed $D_5(K)/D_9(K)$ has dimension
$204+690+2340+8230=11464$, and its homology image has dimension 43.
Its kernel over that image, together with the $7145$ binary dimensions
of the kernel in (D.33.8), has order
$2^{11421+7145}=2^{18566}$, equal to (D.33.10).

Let $S=\beta s\beta^{-1}$ and let $T$ denote conjugation by $t$.
For left corrections $b\beta,aA$ in $\mathcal A$, the exact change of
the cusp error is
\[
                         (1-S)b+(S-T)a.               \tag{D.33.11}
\]
To verify the formula, multiply $bSb^{-1}aRa^{-1}t^{-1}$, where
$R=ArA^{-1}$, and move its four correction factors to the left.
The original error $SRt^{-1}$ lies in $\mathcal A$, so its conjugation
on this abelian group is trivial. This gives (D.33.11). Changing the
multiplier by 144 adds the actual $t^{-144}$ coordinate. Its degree-eight
term equals that of $t^{48}$, and it is retained in every calculation.
No quadratic terms have been removed: they vanish because the chosen
corrections commute in each of the two full target groups.

Conjugation on $U_8$ now retains the degree-eight effect of $N$.
In particular $[N,U_8]$ is the fourth subgroup of dimension 1620.
The program uses the actual degree-two inner factor in
$\beta s\beta^{-1}=ks$, the complete involution on the four $K$
generators, and the actual $t$-conjugate generator words. The resulting
map (D.33.11) on $U_8^2$ has rank 8981 and cokernel dimension 124.
All 124 annihilating functionals are obtained by exact elimination.

For the high $C_1$ words, retain their complete order-four coordinates
in $H^{(2)}$ and their full degree-five through degree-eight expansions
in $T_3$. The parity map on the 9416 lower variables has rank 4687.
The upper variables are twice every original $C_1$ coordinate in both
conjugators and the multiplier-144 variable, hence 9507 variables.
In the combined space consisting of the 4753 upper class-two coordinates
and the 124 $U_8$ obstruction coordinates, their image has rank 4738.
Its cokernel has dimension $15+124=139$.

The exact Bockstein calculation on the lower parity relations has rank
139: the combined lower matrix has rank $4687+139=4826$. The functionals
on $U_8$ are extended linearly to the ambient tensor space. This is
legitimate here because satisfying the full class-two equation puts
the final error in $U_8$, by (D.33.9). Each class-two part of an upper
annihilator evaluates evenly on a high lower column: twice that high
word has zero $T_3$ image. Division by two in this even image is the
explicit Bockstein used by the program. Thus these are simultaneous
equations for the actual order-four group, rather than just equations
for its first binary quotient.

Solving these equations, then the upper equation, then the $U_8$
equation gives the recorded multiplier-one words. The two independent
complete-word evaluations in the statement verify every coefficient,
not merely the annihilating functionals. The induced homology map on
$N/\Phi(N)$ is unchanged and has rank 97. The quotient action on $Q$
is unchanged. Characteristicity of $D_9(K)$ and of the defining
class-two subgroup of $N$ makes the resulting maps endomorphisms of
both stated quotients. The finite two-kernel Frattini argument makes
them automorphisms.

To continue while retaining $T_3$, use D.29's stable two-level module
$C_k\simeq(\mathbb Z/4)^{4753}$. At $k=2$, restrict the lower
corrections to the squares of the 4708 high $C_1$ words. Their degree
is at least ten. Exclude the lower multiplier-144 change. The restricted
parity rank is 4687 on 9416 variables; its relation space has dimension
4729. Its Bockstein onto the full 15-dimensional upper obstruction
space has rank 15. This exact computation uses the unchanged $G_2$
data, on which the stable module action depends. All upper $C_3$
words have degree at least sixteen, as does $t^{288}$.

An actual first continuation gives the stated $18,077,708$-letter
cusp identity over $\mathbb Z/8$. The complete new conjugator
corrections and $t^{288}$ are also expanded through degree eight;
each is the identity in $T_3$. For subsequent stages $k\geq3$, all
lower and upper corrections already lie in $D_{16}(K)$. D.29's
one-level-lookahead construction therefore applies while fixing $T_3$
at every step. The initial $H^{(2)}$ central coordinates can change
once in this procedure; after lookahead each lower level stabilizes.
The displayed degree-eight action remains fixed throughout. The
multiplier increments are divisible by 288, so their two-adic limit
lies in $1+32\mathbb Z_2$.

Taking the inverse limit gives (D.33.3), with the specified convergent
conjugators. Its third and fourth subgroups are their computed images
in $T_3$; the fifth vanishes in both factors. This proves (D.33.4).
The inclusion $D_9(K)\subset D_8(K)$ gives (D.33.5), whose kernel is
the 7245-dimensional pure degree-eight part of $U_8$. Direct expansion
of $bS(b)^{-1}$ for the recorded correction $b$ gives a first nonzero
term in degree five. Later stable corrections vanish there, proving
the stated comparison with D.31 by this explicit quotient map.

Choose the original odd-primary conjugators as separate components.
The new two-primary words need not preserve their own odd-primary
abelianization. The unchanged $G_2$ and separately retained $G_3$
coordinates agree with D.28 on $H\times_QG_3$, the exact maximal
common quotient already proved there. The corresponding fibre product
thus carries one action. The global cusp gauge of D.26 and D.29
produces compatible finite shadows, with their nonarithmetical CM
value retained whenever the quotient retains those coordinates.

Finally the remaining source kernel is explicit and nonzero. Choose
three distinct weight-three generators from (D.33.6). Their nonzero
basic triple commutator belongs to
$\gamma_3N^{(2)}\cap D_9(K^{(2)})$: nonvanishing follows from its
nonzero ordinary free Lie term in that free pro-two basis. Hence
(D.33.3) retains a nontrivial relative pro-two kernel. D.45 and
its source-transport supplement exclude a lift of this recorded
action across the full kernel; another action with the prescribed
CM pair remains in the realization problem.
\end{proof}
